\section{ベイズ的判断とベイズファクター}
\label{b28_bayes_judgement}

\subsection{授業内容}

\subsubsection{概要}
前回JASPの基本的な使用方法を学んだ上で，本コマではベイズ統計学における判断方法について理解を深める。まず，事後分布に基づく確信区間（Credible Interval）と実質的等価性の領域（ROPE: Region of Practical Equivalence）を用いた判断方法について学ぶ。次に，モデル比較の手法としてのベイズファクターの概念と解釈，そして周辺尤度の計算方法としてのサヴェージ・ディッキー法について理解する。特に，従来の帰無仮説検定では帰無仮説を棄却することしかできなかったのに対し，ベイズ法では帰無仮説を積極的に支持することも可能である点を強調する。最後に，情報仮説を含むより柔軟なモデル比較や，実践的なデータ分析におけるベイズ的アプローチの利点について学ぶ。これらをJASPを用いた実例を通じて理解することで，より洗練された統計的判断力を身につける。

\subsubsection{コマ主題細目}
\begin{description}
	\item[確信区間とROPEに基づく判断] ベイズ推定では，パラメータの事後分布から確信区間（Credible Interval）を算出できることを学ぶ。例えば95\%確信区間は，「パラメータが区間内に95\%の確率で存在する」と直接的に解釈できる。これは頻度主義的な信頼区間とは解釈が異なり，より直感的である。また，実質的等価性の領域（ROPE: Region of Practical Equivalence）の概念を導入し，パラメータがゼロや特定の値と「実質的に等価」かどうかを判断する方法を学ぶ。具体的には，95\%確信区間がROPEに完全に含まれる場合は「実質的に等価」，完全に外れる場合は「実質的に異なる」，部分的に重なる場合は「判断保留」と判断できる。これにより，単なる「有意/非有意」の二分法ではなく，パラメータの大きさと不確実性を考慮した判断が可能になることを理解する。JASPを用いて，実際のデータセットに対するこれらの判断方法を実演する。

	\item[ベイズファクターの概念と解釈] ベイズ統計学における仮説検証の方法として，ベイズファクター（Bayes Factor, BF）の概念と解釈について学ぶ。ベイズファクターは，二つのモデルの相対的な証拠の強さを表す比率であり，$BF_{10} = \frac{p(data|H_1)}{p(data|H_0)}$と定義される。ここで$p(data|H)$は，モデル$H$のもとでデータが観測される確率（周辺尤度）である。$BF_{10} > 1$のとき対立仮説$H_1$が，$BF_{10} < 1$のとき帰無仮説$H_0$が支持される。$BF_{10}$の値の大きさは証拠の強さを表し，例えば$BF_{10} = 10$は「対立仮説のための証拠が帰無仮説のための証拠の10倍ある」ことを意味する。ベイズファクターの解釈基準（例：1-3は「わずかな証拠」，3-10は「中程度の証拠」，10以上は「強力な証拠」）について学び，具体的な研究例での解釈方法を理解する。JASPでは，分析結果にベイズファクターが自動的に含まれることを確認し，その読み方を実践的に学ぶ。
	      \begin{flushright}
		      $\to$ \textcite{Lee201709}のP.88--103.
	      \end{flushright}

	\item[周辺尤度とサヴェージ・ディッキー法] ベイズファクターを計算するには，各モデルの周辺尤度$p(data|H)$を求める必要がある。周辺尤度は，パラメータ空間全体にわたる尤度と事前分布の積の積分$p(data|H) = \int p(data|\theta, H)p(\theta|H)d\theta$として定義される。しかし，複雑なモデルではこの積分の計算が困難であるため，近似計算法が必要となる。特に，ネストした（入れ子になった）モデル比較の場合，サヴェージ・ディッキー法（Savage-Dickey Density Ratio Method）が有用である。これは，対立仮説の事前分布と事後分布のパラメータ特定値（例えばゼロ）における比率からベイズファクターを近似する方法である。JASPでは，t検定や分散分析などのベイズ版で，このサヴェージ・ディッキー法が内部的に使用されていることを理解する。実例として，JASPを用いた単純なt検定と分散分析のベイズ版を実行し，出力されるベイズファクターの解釈方法を学ぶ。
	      \begin{flushright}
		      $\to$ \textcite{Maeda2017}のP.342--364，\textcite{Lee201709}のP.104--120.
	      \end{flushright}

	\item[情報仮説と実践的応用] ベイズ統計学の大きな利点として，帰無仮説対対立仮説という二分法だけでなく，より洗練された「情報仮説（informed hypothesis）」の検証が可能である点を学ぶ。例えば，パラメータが特定の範囲内にあるという仮説や，複数のパラメータ間の大小関係に関する仮説など，研究者の事前知識や理論に基づいた具体的な仮説を検証できる。これにより，単なる「効果があるかないか」ではなく，「理論と一致する効果があるか」をより直接的に検証できる。JASPではこのような情報仮説の検証も可能であり，その設定方法と解釈を学ぶ。また，ベイズ統計学の実践的応用として，サンプルサイズが小さい場合の推測，複数の実験結果の統合，事前知識の体系的な組み込みなど，従来の頻度主義的アプローチでは難しかった分析が可能になることを理解する。具体的な研究例を通じて，ベイズ的アプローチが心理学研究にもたらす新たな可能性に言及する。

\end{description}
\subsubsection{キーワード}
\begin{itemize}
	\item 確信区間（Credible Interval）
	\item 実質的等価性の領域（ROPE）
	\item ベイズファクター（Bayes Factor）
	\item サヴェージ・ディッキー法
	\item 情報仮説（Informed Hypothesis）
\end{itemize}
\subsection{授業情報}
\paragraph{コマの展開方法}
講義とJASPを用いた実習
\paragraph{標準シラバスにおける位置づけ}
科目番号5；心理学統計法；(1)心理学で用いられる統計手法；J より高度な記述や推定を目指して
\subsubsection{予習・復習課題}
\paragraph{予習}
前回学んだJASPの基本操作を復習し，スムーズに使えるようにしておくこと。また，ベイズ統計学の基本概念（事前分布，事後分布，ベイズの定理など）についても再確認しておくとよい。特に，従来の帰無仮説検定のロジック（帰無仮説が正しいという仮定のもとで，観測されたデータが得られる確率を評価する）と，ベイズ的アプローチのロジック（データが与えられたときの各モデルの相対的な確からしさを評価する）の違いについて理解を深めておくことが重要である。

\paragraph{復習}
授業で学んだ内容を定着させるために，以下の点について理解を深めておくこと：

\begin{enumerate}
\item 確信区間とROPEを用いた判断方法の手順と解釈を自分の言葉で説明できるようにする
\item ベイズファクターの計算方法，解釈基準，$BF_{10}$と$BF_{01}$の関係を理解する
\item サヴェージ・ディッキー法の基本原理と，なぜこの方法が周辺尤度の計算を簡略化できるかを説明できるようにする
\item 情報仮説の設定方法と，それが従来の帰無仮説/対立仮説の枠組みとどのように異なるかを理解する
\item JASPを用いて実際のデータセットに対するベイズ的判断を実施し，結果を適切に解釈できるようにする
\end{enumerate}

特に，JASP内蔵のサンプルデータセットや，自分で入手した心理学的データを用いて，t検定，分散分析，回帰分析などのベイズ版を実行し，ベイズファクターや確信区間の解釈を練習するとよい。また，同じデータに対して頻度主義的アプローチとベイズ的アプローチの両方を適用し，結果の違いとその理由について考察することで，両者の特徴をより深く理解することができるだろう。ベイズ統計学は心理学研究における新たな標準となりつつあるため，実践的なスキルを身につけることが重要である。